\documentclass[11pt]{article}
% Shared preamble for the bite-sized LaTeX doc series (docs/latex/*.tex).
% \input this at the top of each numbered file rather than duplicating the
% package list -- keeps every file's own content close to (and comfortably
% under) 10 pages.
\usepackage[a4paper,margin=2.5cm]{geometry}
\usepackage{amsmath,amssymb}
\usepackage{hyperref}
\usepackage[backend=biber,style=numeric,sorting=none]{biblatex}
\usepackage{listings}
\usepackage{xcolor}
\usepackage{enumitem}
\addbibresource{references.bib}

\lstdefinelanguage{Rust}{
  keywords={fn,let,mut,pub,struct,enum,impl,match,if,else,for,while,return,use,mod,self,Self,const,static},
  keywordstyle=\color{blue}\bfseries,
  comment=[l]{//},
  commentstyle=\color{gray}\itshape,
  string=[b]",
  stringstyle=\color{teal},
  basicstyle=\ttfamily\small,
  showstringspaces=false,
  breaklines=true,
}
\lstset{language=Rust}

\hypersetup{colorlinks=true, linkcolor=blue, citecolor=blue, urlcolor=blue}


\title{The Non-Analytic Critical-Region Term \\ \large Part 1: Theory}
\author{outram-park-fork-coolprop}
\date{2026-07-10}

\begin{document}
\maketitle

\begin{abstract}
A pure fluid's thermodynamic state is fully determined by its reduced
Helmholtz free energy $\alpha(\delta,\tau) = \alpha_0(\delta,\tau) +
\alpha_r(\delta,\tau)$, where $\delta = \rho/\rho_r$ is the reduced density
and $\tau = T_r/T$ the inverse reduced temperature. Near the critical point,
ordinary polynomial and Gaussian-bell terms in $\alpha_r$ cannot reproduce
the sharp, non-classical curvature of real fluid behaviour. IAPWS-95 and
Span--Wagner-style equations of state add a dedicated \emph{non-analytic}
term for exactly this region. This document derives that term's structure;
Part 2 covers the numerical subtleties of evaluating it, Part 3 the Rust
implementation, and Part 4 walks through the verified result.
\end{abstract}

\section{Why an ordinary polynomial term is not enough}

Away from the critical point, $\alpha_r$ is well approximated by sums of
\begin{align}
  \alpha_r^{\text{power}}(\delta,\tau) &= \sum_i n_i\, \delta^{d_i}\, \tau^{t_i}
    \, e^{-\delta^{l_i}}, \\
  \alpha_r^{\text{Gaussian}}(\delta,\tau) &= \sum_i n_i\, \delta^{d_i}\, \tau^{t_i}
    \, e^{-\eta_i(\delta-\varepsilon_i)^2 - \beta_i(\tau-\gamma_i)^2}.
\end{align}
Both forms are smooth, analytic functions of $(\delta,\tau)$ everywhere.
Real fluids, however, exhibit \emph{non-classical critical scaling}: the
isochoric heat capacity $c_v$ diverges at the critical point with a
power-law exponent that no finite sum of analytic terms can reproduce
exactly. IAPWS-95 (and Span--Wagner-style multiparameter EOS more generally)
address this with a small number of specially-shaped terms concentrated
tightly around $(\delta,\tau) = (1,1)$ — the reduced critical point — chosen
so the \emph{overall fit} closely tracks measured near-critical data without
needing to solve the full non-classical scaling problem exactly
\cite{wagner2002iapws}.

\section{The non-analytic term's functional form}

Each non-analytic term $i$ contributes
\begin{equation}
  \alpha_{r,i} = n_i\, \Delta^{b_i}\, \delta\, \psi
  \label{eq:nonanalytic-term}
\end{equation}
built from three auxiliary functions, each itself a function of $(\delta,\tau)$:
\begin{align}
  \theta &= (1-\tau) + A_i\left[(\delta-1)^2\right]^{1/(2\beta_i)},
  \label{eq:theta} \\
  \Delta &= \theta^2 + B_i\left[(\delta-1)^2\right]^{a_i},
  \label{eq:Delta} \\
  \psi &= \exp\!\left(-C_i(\delta-1)^2 - D_i(\tau-1)^2\right).
  \label{eq:psi}
\end{align}
Reading these three pieces separately clarifies what each is doing:

\begin{itemize}[leftmargin=1.5em]
  \item $\psi$ (Eq.~\ref{eq:psi}) is an ordinary two-dimensional Gaussian
    bell centred exactly on the critical point $(\delta,\tau)=(1,1)$, with
    independent widths $C_i, D_i$ in the density and temperature directions.
    It is what confines the whole term's influence to a small neighbourhood
    of the critical point — away from it, $\psi \to 0$ rapidly and
    $\alpha_{r,i} \to 0$ regardless of the other factors.
  \item $\theta$ (Eq.~\ref{eq:theta}) measures a kind of ``distance along the
    critical isotherm,'' combining how far $\tau$ is from $1$ with a
    density-offset correction raised to the fractional power
    $1/(2\beta_i)$ — this fractional exponent is precisely the
    non-analytic ingredient: no integer power of $(\delta-1)$ can produce it.
  \item $\Delta$ (Eq.~\ref{eq:Delta}) combines $\theta^2$ with a second,
    purely density-dependent fractional-power term, and is itself raised to
    a further fractional power $b_i$ in Eq.~\ref{eq:nonanalytic-term}. This
    double layering of fractional powers is what lets the term reproduce
    the steep, non-classical curvature of $c_v$ and related derivatives near
    the critical point that Eqs.~1--2 cannot.
\end{itemize}

For Water (IAPWS-95), CoolProp's \texttt{Water.json} carries exactly two
such terms ($i=1,2$), with coefficients $A_i \approx 0.32$,
$B_i \approx 0.2$, $\beta_i = 0.3$, $a_i = 3.5$, $b_i \in \{0.85, 0.95\}$,
$C_i \in \{28, 32\}$, $D_i \in \{700, 800\}$ \cite{coolprop}. The large $C_i$,
$D_i$ values confirm the Gaussian $\psi$ factor is indeed very tightly
peaked — its half-width in $\delta$ is only $\sim 1/\sqrt{C_i} \approx 0.19$
reduced-density units.

\section{What this buys, physically}

Without this term, an EOS evaluated \emph{exactly at} the critical point
$(\delta,\tau)=(1,1)$ is missing a contribution that is small in absolute
magnitude everywhere except in a narrow band around the critical point —
but inside that band, it is the difference between a qualitatively wrong
critical-point pressure and one that reproduces the defining IAPWS-95 value
to machine precision (Part 4 shows this numerically: $5.2\times10^{-14}$
relative error once implemented, versus a documented $\sim 10^{-4}$ residual
before). The same non-analytic shape appears (with different coefficients)
in Span \& Wagner's CO$_2$ equation of state \cite{span1996wagner} and other
IAPWS/Span--Wagner-style multiparameter EOS — it is a standard building
block, not a Water-specific hack.

\section{What Part 2 covers}

Equations~\ref{eq:theta}--\ref{eq:psi} are well-defined mathematically at
every $(\delta,\tau) \neq (1,1)$, but the term needs first and second
partial derivatives in $\delta$ and $\tau$ (for pressure, $c_v$, $c_p$, and
the speed of sound) — several of the intermediate quantities in those
derivatives are formally $0/0$ or $0^0$ \emph{exactly at} the critical
point. Part 2 works through the derivative chain and the numerical guard
needed to evaluate it there.

\printbibliography
\end{document}
